\section{本讲小结}
磁场计算首先把电流分布与场点几何分清，再决定采用微元积分、对称环路还是正负分布的叠加。毕奥—萨伐尔定律保留每个电流元的方向，安培环路定理则把有向电流与环量联系起来；只有对称性足够时，环量才直接给出某一点的场。

有限与无限、离散与连续、实心与空心等条件改变的是积分边界或包围电流，不能只换一个长度就沿用结论。圆心、轴线、柱面和无穷远的极限，可以检查系数、方向和近似是否相容。
